%% ============================================================
%% Supplementary Material for:
%% Continuous-Time Bayesian Networks with Structured Shrinkage
%% Priors for Modelling Multimorbidity Trajectories in
%% Large-Scale Electronic Health Records
%% ============================================================
\documentclass[12pt]{article}

\usepackage{amsmath,amssymb,amsthm}
\usepackage{geometry}
\geometry{margin=1in}
\usepackage{graphicx}
\usepackage{float}
\usepackage{caption}
\usepackage{booktabs}
\usepackage{longtable}
\usepackage{array}
\usepackage{xcolor}
\usepackage{hyperref}
\hypersetup{colorlinks=true,linkcolor=blue,citecolor=blue,urlcolor=blue}
\usepackage{natbib}
\usepackage{authblk}
\usepackage{setspace}
\usepackage{listings}

\renewcommand{\thesection}{S\arabic{section}}
\renewcommand{\thesubsection}{S\arabic{section}.\arabic{subsection}}
\renewcommand{\thesubsubsection}{S\arabic{section}.\arabic{subsection}.\arabic{subsubsection}}
\renewcommand{\thetable}{S\arabic{table}}
\renewcommand{\thefigure}{S\arabic{figure}}
\renewcommand{\theequation}{S\arabic{equation}}

\title{\textbf{Supplementary Material:\\
Continuous-Time Bayesian Networks with Structured Shrinkage
Priors for Modelling Multimorbidity Trajectories in Large-Scale
Electronic Health Records}}

\author[1]{O.~R.~Olaniran}
\author[1]{S.~S.~Paria}
\author[2]{M.~Khondoker}
\author[2]{A.~MacGregor}
\author[1]{A.~Lewin}
\affil[1]{Department of Medical Statistics, London School of Hygiene \&
Tropical Medicine, London, UK}
\affil[2]{University of East Anglia, Norwich, UK}
\date{}

\begin{document}

\maketitle
\thispagestyle{empty}

\tableofcontents

\clearpage

%% ============================================================
\section{Detailed Specification of the Four Shrinkage Priors}
\label{sec:supp_priors}
%% ============================================================

This section gives the full hierarchical specification, marginal density, and
closed-form full conditionals for each of the four shrinkage priors compared
in the main text. Throughout, $\beta_j$ denotes a coefficient with associated
interaction order $p_j \in \{1, \dots, P\}$ (with $p_j = 1$ for main effects
and $p_j = r$ for an $r$-way interaction), $\theta > 0$ is the order-penalty
hyperparameter, and $\boldsymbol{\gamma}_m$ denotes the covariate coefficients
which are not subject to order-based penalisation. We write $n_p = |\{j : p_j = p\}|$
for the number of coefficients of order $p$, and $n_m$ for the effective sample
size for condition $m$ (number of at-risk intervals).

\subsection{Hierarchical Structured (HS) Prior}

The hierarchical structured prior places an order-specific inverse-gamma
hyperprior on the slab variance.

\subsubsection{Hierarchy}
\begin{align}
\beta_j \mid \sigma^2_{p_j} &\sim \mathcal{N}\!\left(0,\; \sigma^2_{p_j}\, e^{-\theta p_j}\right), \qquad j = 1, \dots, J,\\
\sigma^2_p &\sim \frac{1}{\Gamma(p+1)}\,\mathrm{Inv\text{-}Gamma}(a_0, b_0), \qquad p = 0, 1, \dots, P,\\
\gamma_k \mid \sigma^2_\gamma &\sim \mathcal{N}(0, \sigma^2_\gamma), \qquad k = 1, \dots, K,\\
\sigma^2_\gamma &\sim \mathrm{Inv\text{-}Gamma}(a_0, b_0).
\end{align}
The factorial scaling $\Gamma(p+1)^{-1}$ on the hyperprior reflects the
combinatorially increasing number of potential $p$-way interactions and
imposes additional penalisation at higher orders beyond the exponential decay
$e^{-\theta p_j}$. Together these two mechanisms ensure that main effects are
only mildly regularised while higher-order interactions are strongly shrunk
towards zero. The hyperparameters are typically set to $a_0 = b_0 = 0.001$ to
concentrate the hyperprior mass near zero.

\subsubsection{Marginal density}
Marginalising over $\sigma^2_{p_j}$ yields a scaled Student-$t$ marginal:
\begin{equation}
p(\beta_j) = \frac{\Gamma(a_0 + \tfrac{1}{2})}{\Gamma(a_0)\sqrt{2\pi b_0\,\Gamma(p_j+1)\,e^{\theta p_j}}} \left(1 + \frac{\beta_j^2}{2b_0\,\Gamma(p_j+1)\,e^{\theta p_j}}\right)^{-(a_0+\frac{1}{2})},
\end{equation}
a $t_{2a_0}$ distribution with scale that shrinks with interaction order. The
heavy-tailed marginal allows individual coefficients to escape shrinkage when
strongly supported by data.

\subsubsection{Full conditionals}
\begin{align}
\sigma^2_p \mid \boldsymbol{\beta} &\sim \mathrm{Inv\text{-}Gamma}\!\left(a_0 + \frac{n_p}{2},\;\; b_0 + \frac{e^{\theta p}}{2} \sum_{j:\,p_j=p} \beta_j^2\right),\\
\sigma^2_\gamma \mid \boldsymbol{\gamma} &\sim \mathrm{Inv\text{-}Gamma}\!\left(a_0 + \frac{K}{2},\;\; b_0 + \frac{1}{2}\sum_{k=1}^K \gamma_k^2\right).
\end{align}
The full conditional for $\beta_j$ is non-conjugate due to the Poisson
likelihood and is sampled by NUTS.

\subsection{Bayesian LASSO}

The Bayesian LASSO \citep{ParkCasella2008} induces $L_1$-type regularisation
through a scale mixture of normals representation.

\subsubsection{Hierarchy}
\begin{align}
\beta_j \mid \tau^2_j &\sim \mathcal{N}\!\left(0,\; \tau^2_j\, e^{-\theta p_j}\right), \qquad j = 1, \dots, J,\\
\tau^2_j \mid \lambda^2_{p_j} &\sim \mathrm{Exp}\!\left(\frac{\lambda^2_{p_j}}{2}\right),\\
\lambda^2_p &\sim \mathrm{Inv\text{-}Gamma}(a_0, b_0), \qquad p = 0, 1, \dots, P,\\
\gamma_k \mid \lambda^2_\gamma &\sim \mathrm{Laplace}(0, \sqrt{\lambda^2_\gamma}),\qquad k = 1, \dots, K,\\
\lambda^2_\gamma &\sim \mathrm{Inv\text{-}Gamma}(a_0, b_0).
\end{align}
The per-order penalty $\lambda^2_p$ controls the overall degree of shrinkage
for coefficients of order $p$.

\subsubsection{Marginal density}
Marginalising over $\tau^2_j$ yields the double-exponential (Laplace) marginal:
\begin{equation}
p(\beta_j \mid \lambda_{p_j}) = \frac{\lambda_{p_j}\, e^{\theta p_j/2}}{2} \exp\!\left(-\lambda_{p_j}\, e^{\theta p_j/2}\, |\beta_j|\right),
\end{equation}
so the effective Laplace scale shrinks exponentially with order.

\subsubsection{Full conditionals}
\begin{align}
\tau^{-2}_j \mid \beta_j, \lambda_{p_j} &\sim \mathrm{Inv\text{-}Gaussian}\!\left(\frac{\lambda_{p_j}\, e^{\theta p_j/2}}{|\beta_j|},\;\; \lambda^2_{p_j}\, e^{\theta p_j}\right),\\
\lambda^2_p \mid \{\tau^2_j : p_j = p\} &\sim \mathrm{Inv\text{-}Gamma}\!\left(a_0 + n_p,\;\; b_0 + \frac{1}{2}\sum_{j:\,p_j=p} \tau^2_j\right),\\
\lambda^2_\gamma \mid \boldsymbol{\gamma} &\sim \mathrm{Inv\text{-}Gamma}\!\left(a_0 + K,\;\; b_0 + \frac{1}{2}\sum_{k=1}^K \gamma_k^2\right).
\end{align}
The shrinkage factor used for the pseudo-PIP construction (Section~3.3 of the
main text) is $\kappa_j = (1 + 2\lambda^2_{p_j}\, e^{-\theta p_j})^{-1}$.

\subsection{Regularised Horseshoe}

The regularised horseshoe \citep{PiironenVehtari2017} employs a global--local
shrinkage structure with a slab component that caps the local scale.

\subsubsection{Hierarchy}
\begin{align}
\beta_j &= z_j\, \tilde{\lambda}_j\, \tau_j, \qquad z_j \sim \mathcal{N}(0, 1),\\
\tilde{\lambda}^2_j &= \frac{c^2 \lambda^2_j}{c^2 + \tau^2_j \lambda^2_j},\\
\tau_j &= \tau_0\, e^{-\theta p_j},\\
\lambda_j &\sim \mathrm{Cauchy}^+(0, 1),\\
\tau_0 &\sim \mathrm{Cauchy}^+(0, \tau_{\mathrm{scale}}),\\
c^2 &\sim \mathrm{Inv\text{-}Gamma}\!\left(\frac{\nu_{\mathrm{slab}}}{2},\; \frac{\nu_{\mathrm{slab}}\, s^2_{\mathrm{slab}}}{2}\right),\\
\gamma_k &\sim \mathcal{N}(0, 1).
\end{align}
The non-centred parameterisation $\beta_j = z_j \tilde{\lambda}_j \tau_j$
improves HMC geometry when $\tau_j$ is small. The slab variance $c^2$ caps the
effective local scale at $\sqrt{c^2}$. Defaults are $\nu_{\mathrm{slab}} = 4$,
$s_{\mathrm{slab}} = 2$, and $\tau_{\mathrm{scale}} = p_0 / J \cdot n_m^{-1/2}$
where $p_0$ is a prior guess at the number of non-zero coefficients.

\subsubsection{Auxiliary representation and full conditionals}
The half-Cauchy distributions on $\lambda_j$ and $\tau_0$ admit an inverse-gamma
scale-mixture representation. Introducing auxiliary variables $\nu_j > 0$ and
$\xi > 0$:
\begin{align}
\lambda^2_j \mid \nu_j &\sim \mathrm{Inv\text{-}Gamma}\!\left(\tfrac{1}{2},\; \tfrac{1}{\nu_j}\right), \quad \nu_j \sim \mathrm{Inv\text{-}Gamma}\!\left(\tfrac{1}{2},\; 1\right),\\
\tau^2_0 \mid \xi &\sim \mathrm{Inv\text{-}Gamma}\!\left(\tfrac{1}{2},\; \tfrac{1}{\xi}\right), \quad \xi \sim \mathrm{Inv\text{-}Gamma}\!\left(\tfrac{1}{2},\; 1\right).
\end{align}
The closed-form full conditionals for the auxiliary variables are
\begin{align}
\lambda^{-2}_j \mid \beta_j, \tau_j, \nu_j &\sim \mathrm{Inv\text{-}Gamma}\!\left(1,\;\; \frac{1}{\nu_j} + \frac{\beta_j^2}{2\tau^2_j}\right),\\
\nu^{-1}_j \mid \lambda_j &\sim \mathrm{Inv\text{-}Gamma}\!\left(1,\;\; 1 + \frac{1}{\lambda^2_j}\right),\\
\tau^{-2}_0 \mid \{\beta_j\}, \{\lambda_j\}, \xi &\sim \mathrm{Inv\text{-}Gamma}\!\left(\frac{J+1}{2},\;\; \frac{1}{\xi} + \frac{1}{2}\sum_{j=1}^J \frac{\beta_j^2}{\lambda^2_j\, e^{-2\theta p_j}}\right),\\
\xi^{-1} \mid \tau_0 &\sim \mathrm{Inv\text{-}Gamma}\!\left(1,\;\; 1 + \frac{1}{\tau^2_0}\right).
\end{align}
The posterior shrinkage factor for the pseudo-PIP construction is
$\kappa_j = (1 + n_m\, \tau^2_j\, \tilde{\lambda}^2_j)^{-1}$.

\subsection{Continuous Spike-and-Slab}

The spike-and-slab prior performs explicit probabilistic variable selection
by mixing a diffuse slab with a concentrated spike. We use a continuous
relaxation to enable HMC sampling.

\subsubsection{Hierarchy}
For the interaction and main-effect coefficients:
\begin{align}
\beta_j \mid w_j, \sigma^2_{p_j} &\sim \mathcal{N}\!\left(0,\; w_j \sigma^2_{p_j} + (1-w_j)\epsilon\right), \qquad j = 1, \dots, J,\\
w_j \mid \pi_{p_j} &\sim \mathrm{Beta}\!\left(1,\; \frac{1}{\pi_{p_j}} - 1\right),\\
\pi_p &= \pi_0\, e^{-\theta p}, \qquad \pi_0 \in (0,1),\\
\sigma^2_p &\sim \mathrm{Inv\text{-}Gamma}(a_0, b_0), \qquad p = 0, 1, \dots, P.
\end{align}
The spike variance is set to $\epsilon = 10^{-4}$ so that $\epsilon \ll \sigma^2_{p_j}$
and the continuous relaxation closely approximates the discrete spike-and-slab.
The Beta prior on $w_j$ ensures $\mathbb{E}[w_j] = \pi_{p_j}$.

For the covariate coefficients, an analogous structure with a fixed inclusion
probability $\pi_0$ is specified:
\begin{align}
\gamma_k \mid v_k, \sigma^2_\gamma &\sim \mathcal{N}\!\left(0,\; v_k \sigma^2_\gamma + (1-v_k)\epsilon\right), \qquad k = 1, \dots, K,\\
v_k \mid \pi_0 &\sim \mathrm{Beta}\!\left(1,\; \frac{1}{\pi_0} - 1\right),\\
\sigma^2_\gamma &\sim \mathrm{Inv\text{-}Gamma}(a_0, b_0).
\end{align}

\subsubsection{Marginal density}
Integrating over $w_j$ and $v_k$ yields marginal mixture distributions:
\begin{align}
p(\beta_j) &= \pi_{p_j}\,\mathcal{N}(\beta_j \mid 0, \sigma^2_{p_j}) + (1-\pi_{p_j})\,\mathcal{N}(\beta_j \mid 0, \epsilon),\\
p(\gamma_k) &= \pi_0\,\mathcal{N}(\gamma_k \mid 0, \sigma^2_\gamma) + (1-\pi_0)\,\mathcal{N}(\gamma_k \mid 0, \epsilon).
\end{align}

\subsubsection{Posterior inclusion probability (PIP)}
The PIP for $\beta_j$ is recovered analytically at each MCMC draw via Bayes'
theorem:
\begin{equation}\label{eq:supp_pip}
\mathrm{PIP}_j = \frac{\pi_{p_j}\,\phi(\beta_j;\, 0, \sigma^2_{p_j})}{\pi_{p_j}\,\phi(\beta_j;\, 0, \sigma^2_{p_j}) + (1-\pi_{p_j})\,\phi(\beta_j;\, 0, \epsilon)},
\end{equation}
where $\phi(\cdot;\, 0, v)$ denotes the $\mathcal{N}(0, v)$ density. The
posterior mean PIP, $\widehat{\mathrm{PIP}}_j = \mathbb{E}[\mathrm{PIP}_j \mid \mathcal{D}]$,
is computed by averaging Equation~\eqref{eq:supp_pip} over MCMC draws. A
directed edge is declared active if $\widehat{\mathrm{PIP}}_j \geq 0.5$.

\subsubsection{Full conditionals}
\begin{align}
\sigma^2_p \mid \boldsymbol{\beta}, \mathbf{w} &\sim \mathrm{Inv\text{-}Gamma}\!\left(a_0 + \frac{n_p}{2},\;\; b_0 + \frac{1}{2}\sum_{j:\,p_j=p} \frac{w_j \beta_j^2}{w_j \sigma^2_p + (1-w_j)\epsilon}\right),\\
\sigma^2_\gamma \mid \boldsymbol{\gamma}, \mathbf{v} &\sim \mathrm{Inv\text{-}Gamma}\!\left(a_0 + \frac{K}{2},\;\; b_0 + \frac{1}{2}\sum_{k=1}^K \frac{v_k \gamma_k^2}{v_k \sigma^2_\gamma + (1-v_k)\epsilon}\right).
\end{align}
The full conditionals for $w_j$ and $v_k$ are Beta distributions updated by
the relative likelihoods of $\beta_j$ and $\gamma_k$ under the spike and slab
components. The full conditional for $\beta_j$ is sampled by NUTS.

\subsection{Common Laplace-Approximation Form of the Full Conditional for $\beta_j$}

Across all four priors the Poisson likelihood enters non-conjugately. A
Laplace approximation around the current iterate yields the same approximate
Gaussian form for the full conditional of $\beta_j$:
\begin{equation}
\log p(\beta_j \mid \mathrm{rest}) \approx -\frac{(\beta_j - \mu_j)^2}{2\sigma^2_j} + \mathrm{const},
\end{equation}
with
\begin{equation}
\sigma^2_j = \left(\frac{1}{v_j} + \sum_{i,s} x^2_{j,i,s}\, T_{i,s}\, e^{\eta_{i,s}}\right)^{-1}, \qquad
\mu_j = \sigma^2_j \sum_{i,s} x_{j,i,s}\!\left(N^m_{i,s} - T_{i,s}\, e^{\eta^{(-j)}_{i,s}}\right),
\end{equation}
where $\eta^{(-j)}_{i,s}$ is the linear predictor excluding the contribution of
$\beta_j$, and the effective prior variance $v_j$ differs by prior:
\[
v_j = \begin{cases}
\sigma^2_{p_j}\, e^{-\theta p_j} & \text{(hierarchical structured and spike-and-slab)}, \\
\tau^2_j\, e^{-\theta p_j} & \text{(Bayesian LASSO)}, \\
\tilde{\lambda}^2_j\, \tau^2_j & \text{(regularised horseshoe)}.
\end{cases}
\]
This common form makes clear that the four priors differ only in how $v_j$
is determined from auxiliary scale parameters. Exact posterior sampling is
performed via NUTS in Stan.

%% ============================================================
\section{Stan Code for All Four Shrinkage Priors}
\label{sec:supp_stan}
%% ============================================================

The following R/Stan code implements the four shrinkage priors via the
\texttt{rstan} package. Full reproducible code, including the data-preprocessing
pipeline and the post-processing scripts that compute pseudo-PIPs, PIPs, and
evaluation metrics, is available at
\url{https://github.com/InflAim/CTBN-Multimorbidity}.

\begin{footnotesize}
\begin{verbatim}
# =============================================================================
# CTBN BAYESIAN POSTERIOR ESTIMATION  (rstan)
# Supports four prior specifications: spike_slab, structured, lasso, horseshoe
# =============================================================================

# ---- 1. Spike-and-Slab (marginalised normal mixture) ------------------------
.STAN_SPIKE_SLAB <- "
data {
  int<lower=1> N_rows;  int<lower=1> Q;  int<lower=1> P;
  array[N_rows] int<lower=0> N_obs;
  matrix[N_rows, Q] X;  matrix[N_rows, P] Z;
  vector<lower=0>[N_rows] T;
  array[Q] int<lower=0> beta_order;
  real<lower=0> theta;  real<lower=0> v0;  real<lower=0> v1;
  real<lower=0,upper=1> pi_inc;  int<lower=1> max_order;
}
parameters { vector[Q] beta;  vector[P] gamma;  real<lower=0> sigma2; }
model {
  sigma2 ~ inv_gamma(1, 1);
  for (j in 1:Q)
    target += log_sum_exp(
      log1m(pi_inc) + normal_lpdf(beta[j] | 0, sqrt(v0 * sigma2)),
      log(pi_inc)   + normal_lpdf(beta[j] | 0,
                        sqrt(v1 * sigma2 * exp(-theta * beta_order[j]))));
  gamma ~ normal(0, 1);
  target += poisson_log_lpmf(N_obs | X * beta + Z * gamma + log(T));
}
"

# ---- 2. Hierarchical Structured ---------------------------------------------
.STAN_STRUCTURED <- "
parameters {
  vector[Q] beta;  vector[P] gamma;
  vector<lower=0>[max_order+1] sigma2;  real<lower=0> sigma2_gamma;
}
model {
  sigma2 ~ inv_gamma(0.001, 0.001);
  for (j in 1:Q)
    beta[j] ~ normal(0, sqrt(sigma2[beta_order[j]+1]
                             * exp(-theta * beta_order[j])));
  gamma ~ normal(0, sqrt(sigma2_gamma));
  target += poisson_log_lpmf(N_obs | X * beta + Z * gamma + log(T));
}
"

# ---- 3. Bayesian LASSO -------------------------------------------------------
.STAN_LASSO <- "
parameters {
  vector[Q] beta;  vector[P] gamma;
  vector<lower=0>[max_order+1] lambda2;  real<lower=0> lambda2_gamma;
}
model {
  lambda2 ~ inv_gamma(0.001, 0.001);
  for (j in 1:Q)
    beta[j] ~ double_exponential(0,
                sqrt(lambda2[beta_order[j]+1] * exp(-theta * beta_order[j])));
  gamma ~ double_exponential(0, sqrt(lambda2_gamma));
  target += poisson_log_lpmf(N_obs | X * beta + Z * gamma + log(T));
}
"

# ---- 4. Regularised Horseshoe (Piironen & Vehtari 2017) ---------------------
.STAN_HORSESHOE <- "
parameters {
  vector[Q] z_beta;  vector[P] gamma;
  vector<lower=0>[Q] lambda;  real<lower=0> tau;  real<lower=0> c2;
}
transformed parameters {
  vector[Q] beta;
  for (j in 1:Q) {
    real tau_j     = tau * exp(-theta * beta_order[j]);
    real l2        = square(lambda[j]);
    real lt        = c2 * l2 / (c2 + square(tau_j) * l2);
    beta[j]        = z_beta[j] * sqrt(lt) * tau_j;
  }
}
model {
  tau    ~ cauchy(0, tau0);
  c2     ~ inv_gamma(0.5 * slab_df, 0.5 * slab_df * square(slab_scale));
  lambda ~ cauchy(0, 1);
  z_beta ~ std_normal();
  gamma  ~ normal(0, 1);
  target += poisson_log_lpmf(N_obs | X * beta + Z * gamma + log(T));
}
"
\end{verbatim}
\end{footnotesize}

%% ============================================================
\section{Detailed Prior Ranking Discussion}
\label{sec:supp_ranking}
%% ============================================================

This section gives the detailed per-prior commentary supporting the consolidated
ranking summary in the main text (Figure~6 of the main paper).

\paragraph{Spike-and-slab.}
The spike-and-slab ranks first across all three analysis models on RMSE,
coverage distance from $0.95$, FPR, and Poisson log-likelihood, demonstrating
consistent superiority on parameter recovery, calibration, false-discovery
control, and predictive accuracy simultaneously. Its sole and persistent
weakness is TPR, on which it ranks last (4th) across all analysis models,
reflecting the conservative nature of its hard-thresholding selection
mechanism. Brier score ranking is 4th under A1, improving to 3rd under A2 and
2nd under A3. Selection AUC ranks 3rd under A1 and A2, improving to 2nd under
A3. Overall, the spike-and-slab presents the most balanced profile across
estimation and selection metrics, with the trade-off of reduced sensitivity.

\paragraph{Hierarchical Structured.}
The structured normal ranks 2nd on RMSE across all analysis models, and
achieves 1st on selection AUC under A2, indicating competitive PIP ranking
ability under correct specification. However, it ranks poorly on coverage
distance (2.5th under A1, 3rd under A2 and A3) and last on Brier score under
A2 and A3 (4th), suggesting that while it recovers point estimates reasonably
well, its uncertainty quantification and predictive calibration are weaker
than the spike-and-slab. FPR ranks 3rd consistently, indicating moderate
false-discovery control. The structured normal is thus a reasonable choice
when point estimation is the priority but is not recommended when calibrated
uncertainty or predictive performance is critical.

\paragraph{Bayesian LASSO.}
The Bayesian LASSO ranks last (4th) on RMSE across all analysis models and
last on FPR, confirming poor parameter recovery and the weakest
false-discovery control of all priors. Its ranks on TPR (2nd across all
analysis models) and selection AUC (1st under A1 and A2, 3rd under A3) reflect
its tendency to include most parameters, inflating both true and false
positives simultaneously. Brier score is 1st under A1 and A2, suggesting
competitive short-term predictive accuracy despite poor estimation, though
this advantage disappears under A3 where its Brier rank falls to 3rd. The
Bayesian LASSO is not recommended when parameter interpretability or network
sparsity is required.

\paragraph{Regularised horseshoe.}
The regularised horseshoe ranks worst (4th) on coverage distance and selection
AUC across all analysis models, consistent with its near-random pseudo-PIP
discrimination documented in the main-text selection figures. It ranks 2nd
on TPR across all analysis models, but as established earlier this reflects
indiscriminate inclusion rather than genuine sensitivity. A notable exception
is under A3, where the regularised horseshoe ranks 1st on coverage distance
and 1st on Brier score, suggesting that its global shrinkage mechanism
confers some advantage in uncertainty quantification and prediction when the
analysis model is over-specified, though this does not compensate for its
poor selection behaviour.

\paragraph{Overall recommendation.}
No single prior dominates uniformly across all seven metrics. However, the
spike-and-slab presents the most coherent profile for the goals of network
structure learning: it ranks first on RMSE, FPR, coverage calibration, and
Poisson log-likelihood across all or most analysis models, at the cost of
reduced sensitivity (TPR rank 4). The structured normal is a viable secondary
choice when computational considerations favour a continuous prior, provided
the analysis does not require hard variable selection. The Bayesian LASSO and
regularised horseshoe are not recommended for this application given their
systematic failure on FPR and selection AUC respectively.

%% ============================================================
\section{Additional Predictive Performance Details}
\label{sec:supp_pred}
%% ============================================================

The under-specified analysis model A1 produces the largest predictive
degradation: under A1, all priors yield Poisson log-likelihood of $-0.5746$
and Brier scores of $0.1581$--$0.1582$, modestly worse than A2, reflecting the
omission of true interaction terms from the linear predictor. Notably, under
the over-specified analysis model A3, predictive performance is marginally
better than or equivalent to A2 for most priors, with the spike-and-slab
achieving the best Poisson log-likelihood ($-0.5721$) and joint-best Brier
score ($0.1572$) across all analysis-model/prior combinations. The Bayesian
LASSO is the sole exception, showing a deterioration under A3 relative to A2
(Poisson LL $-0.5744$ vs.\ $-0.5731$), consistent with its tendency to
over-shrink or over-fit under an expanded parameter space.

The oracle efficiency analysis (Figure~A1 in the main appendix) confirms
that, across all three analysis models, the Brier-ratio distributions for all
four priors are tightly concentrated at or just above $1.0$ for most
conditions, indicating that the fitted models approach oracle-level predictive
efficiency regardless of prior choice or model specification. Conditions with
greater baseline incidence heterogeneity across replicates (notably DM, DVL,
HL, and OA) exhibit slightly wider Brier-ratio distributions, reflecting
greater Monte Carlo variability in the oracle benchmark rather than genuine
inefficiency of the fitted models. Under A3, the distributions remain tightly
clustered near $1.0$ for the spike-and-slab, structured normal, and
regularised horseshoe, while the Bayesian LASSO shows a marginally wider
spread in several conditions (e.g.\ DE, DVL). No prior exhibits systematic
left-tail mass below $0.9$ in any condition or analysis model, confirming
that none of the fitted models suffer from meaningful predictive inefficiency
relative to the oracle.

%% ============================================================
\bibliographystyle{apalike}
\bibliography{reference}
%% ============================================================

\end{document}
